\documentclass[11pt]{article}

% Change "review" to "final" to generate the final (sometimes called camera-ready) version.
% Change to "preprint" to generate a non-anonymous version with page numbers.
\usepackage{acl}
%\usepackage[preprint]{acl}
% Standard package includes
\usepackage{times}
\usepackage{latexsym}
\usepackage{amsmath}
\usepackage{bbm}
\hyphenation{Qwen SemEval bi-encoder}

% For proper rendering and hyphenation of words containing Latin characters (including in bib files)
\usepackage[T1]{fontenc}

% This assumes your files are encoded as UTF8
\usepackage[utf8]{inputenc}

% This is not strictly necessary, and may be commented out,
% but it will improve the layout of the manuscript,
% and will typically save some space.
\usepackage{microtype}

% This is also not strictly necessary, and may be commented out.
% However, it will improve the aesthetics of text in
% the typewriter font.
\usepackage{inconsolata}

%Including images in your LaTeX document requires adding
%additional package(s)
\usepackage{graphicx}

% If the title and author information does not fit in the area allocated, uncomment the following
%
%\setlength\titlebox{<dim>}
%
% and set <dim> to something 5cm or larger.

\title{VerbaNexAI at SemEval-2026 Task 4: Two-Stage Narrative Similarity via Fine-Tuned Bi-Encoder with MLP Ensemble}

\author{
  Pablo A. Pertuz-Duran, Jairo E. Serrano, \\
  \textbf{Juan C. Martinez-Santos, Edwin Puertas}\\
  Universidad Tecnologica de Bolivar, Cartagena, Colombia \\
  \texttt{(ppertuz,jserrano,jcmartinezs,epuerta)@utb.edu.co}
}

\begin{document}

\maketitle

\begin{abstract}
This paper describes VerbaNex AI's participation in SemEval-2026 Task 4: Narrative Similarity, a shared task on assessing semantic relatedness between short narrative texts. The task comprises two tracks: Track~A requires selecting which of two candidate stories is more similar to an anchor, and Track~B requires producing fixed-size story embeddings whose cosine similarity reflects narrative relatedness. We propose a unified two-stage system built on Qwen3-Embedding-0.6B. The first stage fine-tunes the encoder as a bi-encoder with a 512-dimensional projection head using a composite loss combining margin ranking, pairwise softmax, and multiple negatives ranking objectives. The second stage trains a lightweight MLP head over frozen bi-encoder embeddings using pairwise interaction features, with $k$-fold cross-validation and logit-averaging ensemble inference. The system was trained exclusively on the official supervised data without leveraging the additional 1,900 synthetic triples generated by LLM released by the organizers. 
Although the system ranked first on both tracks in the development phase, its performance did not transfer to the official test set, where it ranked 47 on Track A and 22 on Track B.
\end{abstract}


\section{Introduction}

Assessing whether two narratives convey similar underlying stories is a challenging problem that extends beyond surface-level lexical overlap. Narrative similarity requires modeling events, temporal progression, and causal structure \cite{chambers2008unsupervised}. Unlike sentence-level semantic textual similarity (STS) \cite{cer2017semeval}, narrative relatedness involves reasoning over multi-sentence texts where semantic alignment may be implicit, structurally varied, or abstract.

SemEval-2026 Task 4, Narrative Similarity, addresses this challenge by evaluating systems on their ability to assess semantic relatedness between short narrative texts. The shared task consists of two complementary tracks. Track A formulates the problem as comparative ranking: given an anchor story and two candidate stories, systems must determine which candidate is more similar to the anchor. Track B requires generating fixed-size story embeddings such that cosine similarity reflects narrative relatedness. Together, these tracks evaluate both decision-level ranking performance and the quality of learned semantic representations.

Recent advances in contrastive learning have significantly improved embedding-based similarity modeling \cite{gao2021simcse}. Motivated by this progress, we develop a unified two-stage system built on \texttt{Qwen3-Embedding-0.6B}. In the first stage, we train a bi-encoder to learn narrative-level representations using contrastive objectives tailored to the task. In the second stage, we refine similarity predictions through a lightweight interaction model operating over frozen embeddings. This design allows us to balance representation quality with pairwise decision modeling while maintaining computational efficiency.

The code is publicly available at \url{https://github.com/VerbaNexAI}.


\section{Background}

\subsection{Task and Dataset}

SemEval-2026 Task 4 provides a dataset of short narrative texts annotated for pairwise similarity. The Track~A development set contains triples of the form (anchor, candidate~A, candidate~B), each labeled with a binary indicator of which candidate is narratively closer to the anchor. The Track~B development set contains individual story texts to be embedded into a shared vector space. In addition to the labeled development data, the organizers released a synthetic training set of 1,900 triples generated using large language models, intended to supplement the small development set. Our system was trained exclusively on the development data and did not incorporate the synthetic triples.
% TODO: Cite the task organizers' paper once available

% TODO: Fill in exact dataset statistics from JSONL files
\begin{table}[ht]
  \centering
  \begin{tabular}{lc}
    \hline
    \textbf{Statistic} & \textbf{Value} \\
    \hline
    Dev triples (Track A)           & 200         \\
    Dev stories (Track B)           & 479         \\
    Synthetic triples (LLM-gen.)    & 1,900       \\
    Avg.\ story length (tokens)     & 122.0       \\
    \hline
  \end{tabular}
  \caption{Dataset statistics for the Narrative Similarity task.}
  \label{tab:dataset_stats}
\end{table}

Table~\ref{tab:dataset_stats} summarizes the dataset statistics. The development set is small (200 labeled triples), which motivated our design choices toward parameter-efficient architectures and $k$-fold cross-validation. The availability of 1,900 synthetic triples, which we did not use, represents a roughly 9.5$\times$ increase in training data that could have significantly improved generalization.


\subsection{Related Work}

Sentence-level semantic similarity has been widely studied through benchmarks such as STS Benchmark \cite{cer-etal-2017} and SICK \cite{marelli-etal-2014}. The dominant paradigm uses bi-encoder architectures trained with contrastive or ranking losses, as popularized by Sentence-BERT \cite{reimers-gurevych-2019} and subsequent models such as SimCSE \cite{gao-etal-2021} and E5 \cite{wang-etal-2022-text}. More recent embedding models, including BGE \cite{xiao-etal-2023-bge} and Qwen3-Embedding \cite{qwen3-embedding}, have advanced the state of the art on MTEB benchmarks \cite{muennighoff-etal-2023-mteb} by leveraging larger transformer backbones and improved training objectives.

However, narrative similarity extends beyond sentence-level semantics. Prior work in computational narrative understanding has explored event sequence alignment \cite{chambers-jurafsky-2008} and narrative structure extraction \cite{reagan-etal-2016}. These approaches typically require explicit narrative parsing, making them difficult to integrate into an end-to-end embedding framework.

Our approach builds directly on the bi-encoder paradigm established by Sentence-BERT \cite{reimers-gurevych-2019}, extending it with a task-specific composite loss and a projection head for dimensionality reduction. Unlike prior narrative analysis methods that require explicit structural representations, our system relies entirely on learned embeddings from a general-purpose encoder fine-tuned with narrative similarity annotations. Our initial experiments used BGE-Large \cite{xiao-etal-2023-bge} as the backbone, but upgrading to Qwen3-Embedding-0.6B \cite{qwen3-embedding} yielded substantial gains on the development set. As our test results show, however, these development-set gains did not transfer to the held-out evaluation data, suggesting that the improvements were confounded by overfitting to the small development set.


\subsection{Resources Used}

Our system relies on the following publicly available resources: the Qwen3-Embedding-0.6B pretrained model \cite{qwen3-embedding} served as the backbone encoder; the Sentence-Transformers library \cite{reimers-gurevych-2019} provided the bi-encoder training framework; and PyTorch was used for all custom training loops, including the MLP head and the composite loss implementation. Data augmentation was performed using a T5-based paraphrasing model (Vamsi/T5\_Paraphrase\_Paws) and synonym replacement to expand the limited training set. No external labeled datasets were used; notably, the 1,900 synthetic training triples provided by the organizers were also not incorporated.


\section{System Overview}

Our system addresses both tracks through a unified pipeline: a shared bi-encoder backbone produces story embeddings (Track~B), and a lightweight MLP head operates on these embeddings for pairwise ranking (Track~A). Figure~\ref{fig:pipeline} illustrates the complete architecture.

% TODO: Replace placeholder with actual pipeline diagram
\begin{figure*}[t!]
  \centering
  \fbox{\parbox{0.9\textwidth}{\centering\vspace{2cm}[Pipeline diagram placeholder]\vspace{2cm}}}
  \caption{System pipeline. Left: Track~B bi-encoder with projection head produces 512-dim story embeddings. Right: Track~A MLP head computes pairwise interaction features from the frozen embeddings and outputs a similarity score.}
  \label{fig:pipeline}
\end{figure*}


\subsection{Track B: Bi-Encoder with Projection Head}

The Track~B component is a fine-tuned bi-encoder that maps each story to a fixed-size embedding vector. We use Qwen3-Embedding-0.6B \cite{qwen3-embedding} as the backbone encoder, selected for its strong performance on MTEB benchmarks and its dedicated embedding architecture with last-token pooling. This model was chosen after earlier experiments with BGE-Large-en-v1.5 \cite{xiao-etal-2023-bge} reached a development-set performance plateau at $\sim$0.93 accuracy, suggesting that backbone capacity was the limiting factor.

The model is augmented with a linear projection head that reduces the native encoder output from 1024 dimensions to 512 dimensions, followed by L2 normalization:

\begin{equation}
\mathbf{e} = \text{L2Norm}\bigl(\mathbf{W}_{\text{proj}} \cdot \text{LastTok}(\text{Enc}(\mathbf{x})) + \mathbf{b}\bigr)
\end{equation}

\noindent where $\mathbf{x}$ is the input story text, $\text{LastTok}(\cdot)$ extracts the last-token hidden state following the Qwen3-Embedding design and official usage examples \cite{qwen3-embedding}, and $\mathbf{W}_{\text{proj}} \in \mathbb{R}^{512 \times 1024}$ is the projection matrix. The dimensionality reduction to 512 was chosen to improve generalization on the small training set while maintaining representational capacity.

The bi-encoder is trained with a composite loss combining three objectives:

\begin{equation}
\mathcal{L} = w_m \cdot \mathcal{L}_{\text{margin}} + w_p \cdot \mathcal{L}_{\text{pairwise}} + w_n \cdot \mathcal{L}_{\text{MNR}}
\end{equation}

\noindent where $\mathcal{L}_{\text{margin}}$ is a margin ranking loss (margin = 0.1) that enforces a minimum distance between positive and negative pairs, $\mathcal{L}_{\text{pairwise}}$ is a triple-wise pairwise softmax loss with temperature scaling ($\tau = 0.5$), and $\mathcal{L}_{\text{MNR}}$ is the multiple negatives ranking loss \cite{henderson-etal-2017} that treats all in-batch non-matching pairs as negatives. The weights $w_m = 1.0$, $w_p = 0.5$, and $w_n = 0.5$ were selected empirically. This simplified composite loss was adopted after experiments showed that adding hard-negative mining, SimCSE-style self-supervision, and curriculum learning to the loss did not improve---and in some cases degraded---development performance. At inference, each story is encoded independently and only once; cosine similarity between embedding pairs is used for evaluation.

\paragraph{Data Augmentation.} To mitigate the limited training data, we applied two augmentation strategies: T5-based paraphrasing (generating 2 augmented versions per example) and synonym replacement (replacing up to 3 words per story with probability 0.15). Augmented samples were filtered by cosine similarity to the original (retaining only those in the 0.3--0.9 range) to avoid trivial duplicates or semantically drifted variants.


\subsection{Track A: MLP Head with Pairwise Features}

The Track~A component is a lightweight MLP head that operates on frozen Track~B embeddings to make pairwise ranking decisions. Given a triple (anchor, candidate~A, candidate~B), the system first encodes all three stories independently using the frozen bi-encoder. For each anchor--candidate pair, we construct a pairwise feature vector:

\begin{equation}
\boldsymbol{\phi}(\mathbf{e}_a, \mathbf{e}_c) = [\mathbf{e}_a;\, \mathbf{e}_c;\, |\mathbf{e}_a - \mathbf{e}_c|;\, \mathbf{e}_a \odot \mathbf{e}_c]
\end{equation}

\noindent where $\mathbf{e}_a$ and $\mathbf{e}_c$ are the anchor and candidate embeddings, $|\cdot|$ denotes element-wise absolute difference, $\odot$ denotes element-wise product, and $[;\,]$ denotes concatenation. This four-component representation follows the pairwise interaction features used in natural language inference models \cite{conneau-etal-2017-infersent}, yielding a feature vector of dimension $4d = 2048$ for $d = 512$. The feature vector is processed by a two-layer MLP:

\begin{equation}
\text{MLP}(\boldsymbol{\phi}) = \mathbf{W}_2 \bigl(\text{Dropout}(\text{GELU}(\mathbf{W}_1 \boldsymbol{\phi} + \mathbf{b}_1))\bigr) + \mathbf{b}_2
\end{equation}

\noindent where $\mathbf{W}_1 \in \mathbb{R}^{512 \times 2048}$, $\mathbf{W}_2 \in \mathbb{R}^{1 \times 512}$, and the hidden dimension is 512. For each triple, the MLP produces scores $s_A = \text{MLP}(\boldsymbol{\phi}(\mathbf{e}_{\text{anchor}}, \mathbf{e}_A))$ and $s_B = \text{MLP}(\boldsymbol{\phi}(\mathbf{e}_{\text{anchor}}, \mathbf{e}_B))$. The prediction is $\hat{y} = \mathbbm{1}[s_A > s_B]$.

The MLP is trained with cross-entropy loss over the stacked scores $[s_A, s_B]$ divided by a temperature $\tau = 0.7$, augmented with a KL divergence distillation term from Track~B cosine similarities:

\begin{equation}
\mathcal{L}_A = \mathcal{L}_{\text{CE}} + \lambda \cdot \text{KL}\bigl(\sigma(\mathbf{s} / \tau_d) \,\|\, \sigma(\mathbf{c} / \tau_d)\bigr)
\end{equation}

\noindent where $\mathbf{c}$ are the Track~B cosine similarity scores acting as teacher, $\tau_d = 0.7$ is the distillation temperature, and $\lambda = 0.3$. This encourages the MLP head to remain consistent with the bi-encoder's similarity judgments while learning to correct its errors.

To ensure robustness on the small dataset, we train 5 independent MLP heads using $k$-fold cross-validation. At inference, the scores from all heads are averaged before the final prediction:

\begin{equation}
\hat{y} = \mathbbm{1}\Bigl[\frac{1}{K}\sum_{k=1}^{K} s_A^{(k)} > \frac{1}{K}\sum_{k=1}^{K} s_B^{(k)}\Bigr]
\end{equation}


\subsection{Model Variants Explored}

During development, we explored several alternative configurations before converging on the final system. Table~\ref{tab:variants} summarizes the key variants and their outcomes on the development set. The most significant improvement came from upgrading the backbone from BGE-Large to Qwen3-Embedding-0.6B, while additional training complexity provided marginal or negative returns. However, as described in Section~\ref{sec:results}, these development-set gains did not generalize to the official test set.

\begin{table}[ht]
  \centering
  \small
  \begin{tabular}{lc}
    \hline
    \textbf{Variant} & \textbf{Dev Acc} \\
    \hline
    BGE-Large backbone              & $\sim$0.93 \\
    + Hard-negative mining          & No improvement \\
    + SimCSE self-supervision       & Slight degradation \\
    + Curriculum learning           & No improvement \\
    Qwen3-0.6B backbone             & $\sim$0.94--0.95 \\
    + Simplified composite loss     & Best dev result \\
    \hline
  \end{tabular}
  \caption{Model variants explored during development. All results are on the development set.}
  \label{tab:variants}
\end{table}


\section{Experimental Setup}
\label{sec:experimental_setup}

\subsection{Data Splits and Usage}

The task organizers provided two data sources: a labeled development set (on the order of 100 triples) and a synthetic training set of 1,900 triples generated using large language models. Our system was trained exclusively on the development data; the synthetic training triples were not used. For Track~B, the bi-encoder was trained on the full development set with periodic evaluation on Track~A triple accuracy as a proxy metric. For Track~A, we applied 5-fold stratified cross-validation over the development triples: in each fold, 80\% of the triples were used for training and 20\% for validation. All 5 resulting MLP heads were retained for ensemble inference. For the final test submission, the models trained on the development set were applied directly to the test data without any additional fine-tuning or retraining.

\subsection{Preprocessing}

Input stories were tokenized using the Qwen3-Embedding tokenizer with a maximum sequence length of 384 tokens. Stories exceeding this limit were truncated from the right. No additional text cleaning or normalization was applied beyond standard tokenization, as the narrative texts were already well-formed. For data augmentation, T5-based paraphrasing and synonym replacement were applied to expand the training set by approximately 3$\times$, with cosine similarity filtering to discard low-quality augmentations.

\subsection{Training Configuration}

All experiments were conducted on Google Colab using a single NVIDIA Tesla T4 GPU with 16~GB of VRAM. Training was performed with mixed-precision (fp16) using PyTorch's automatic mixed precision. Gradient checkpointing was enabled to reduce memory consumption during Track~B fine-tuning.

\paragraph{Track B.}
The bi-encoder was fine-tuned using the AdamW optimizer \cite{loshchilov-hutter-2019} with a learning rate of $2 \times 10^{-5}$, weight decay of 0.01 for the backbone and 0.05 for the projection head. We used a batch size of 4 with gradient accumulation over 2 steps (effective batch size of 8), a warmup of 500 steps, and trained for 4 epochs. Gradient clipping was applied with a maximum norm of 1.0.

\paragraph{Track A.}
Each MLP head was trained using AdamW with a learning rate of $1 \times 10^{-4}$, weight decay of 0.01, a warmup ratio of 0.1, and a maximum of 10 epochs with early stopping (patience of 3 epochs, monitored on validation accuracy). The batch size was 8. The temperature for score scaling was set to 0.7 and the distillation weight to 0.3.

Table~\ref{tab:hyperparams} summarizes the hyperparameter configuration.

\begin{table}[ht]
  \centering
  \begin{tabular}{lcc}
    \hline
    \textbf{Hyperparameter} & \textbf{Track B} & \textbf{Track A} \\
    \hline
    Backbone             & Qwen3-0.6B   & Frozen Track B \\
    Embedding dim        & 512          & 512            \\
    Batch size           & 4 ($\times$2 accum) & 8    \\
    Learning rate        & 2e-5         & 1e-4           \\
    Weight decay         & 0.01 / 0.05  & 0.01           \\
    Epochs               & 4            & 10             \\
    Temperature          & 0.5          & 0.7            \\
    Dropout (MLP)        & ---          & 0.2            \\
    Mixed precision      & fp16         & fp16           \\
    Gradient clip        & 1.0          & 1.0            \\
    Early stopping       & ---          & patience = 3   \\
    \hline
  \end{tabular}
  \caption{Hyperparameter configuration for Track~B and Track~A.}
  \label{tab:hyperparams}
\end{table}


\subsection{External Tools and Libraries}

The system was implemented in Python 3.11 (Google Colab environment) using PyTorch 2.x for all model training and inference. The Sentence-Transformers library \cite{reimers-gurevych-2019} was used for the bi-encoder training framework. The Hugging Face Transformers library provided the Qwen3-Embedding model and tokenizer. Data augmentation used the T5\_Paraphrase\_Paws model from the Hugging Face Hub. NumPy was used for embedding I/O (Track~B outputs are stored as \texttt{.npy} arrays) and scikit-learn for cross-validation splits.


\section{Results}
\label{sec:results}

\subsection{Development Set Results}

Table~\ref{tab:dev_results} reports the accuracy on the development set. Track~B achieves a best accuracy of 0.935 using direct cosine similarity between embeddings. Track~A, which adds the MLP head ensemble on top of Track~B embeddings, reaches a mean cross-validation accuracy of 0.945.

\begin{table}[ht]
  \centering
  \begin{tabular}{lc}
    \hline
    \textbf{System} & \textbf{Dev Accuracy} \\
    \hline
    Track B (cosine sim, best)     & 0.935   \\
    Track A (MLP, single fold)     & 0.975 (best fold)   \\
    Track A (MLP, 5-fold CV mean)  & 0.945   \\
    \hline
  \end{tabular}
  \caption{Accuracy on the development set.}
  \label{tab:dev_results}
\end{table}

The $k$-fold cross-validation results for the Track~A MLP head are shown in Table~\ref{tab:kfold}. The fold scores are consistently high (0.875--0.975), with a mean of 0.945 and standard deviation of 0.0367. Despite this apparently stable development performance, the large test-set drop indicates that the model still overfit to the limited development distribution.

\begin{table}[ht]
  \centering
  \begin{tabular}{lc}
    \hline
    \textbf{Fold} & \textbf{Val Accuracy} \\
    \hline
    Fold 0 & 0.950  \\
    Fold 1 & 0.875  \\
    Fold 2 & 0.950  \\
    Fold 3 & 0.975  \\
    Fold 4 & 0.975  \\
    \hline
    \textbf{Mean $\pm$ Std} & \textbf{0.945 $\pm$ 0.0367} \\
    \hline
  \end{tabular}
  \caption{$K$-fold cross-validation results for Track~A MLP head.}
  \label{tab:kfold}
\end{table}

% TODO: Add k-fold accuracy bar chart figure here
% \begin{figure}[ht]
%   \centering
%   \includegraphics[width=\columnwidth]{assets/kfold_accuracy.png}
%   \caption{Per-fold validation accuracy for the Track~A MLP head ensemble.}
%   \label{fig:kfold}
% \end{figure}


\subsection{Official Test Results}

Table~\ref{tab:official} reports the official test set results. The system achieved 53.50\% accuracy on Track~A (rank 47 out of 47) and 59.25\% accuracy on Track~B (rank 22 out of 28). The Track~A result is near random chance for a binary choice task, representing a catastrophic drop from the $\sim$0.95 development accuracy. The Track~B result, while below the median, shows slightly better generalization than Track~A, suggesting that the bi-encoder embeddings retained some useful signal despite overfitting.

\begin{table}[ht]
  \centering
  \begin{tabular}{lccc}
    \hline
    \textbf{Track} & \textbf{Test Acc} & \textbf{Dev Acc} & \textbf{Rank} \\
    \hline
    Track A & 53.50\% & 94.50\% & 47 / 47 \\
    Track B & 59.25\% & 93.50\% & 22 / 28 \\
    \hline
  \end{tabular}
  \caption{Official competition results compared to development accuracy.}
  \label{tab:official}
\end{table}

The contrast between development and test performance (a drop of over 40 percentage points on Track~A) reveals that our system severely overfit to the small development set. We identify two primary causes. First, the system was trained exclusively on the development data (200 triples), which is insufficient to learn generalizable narrative similarity patterns; the 1,900 synthetic training triples provided by the organizers were not used, representing a missed opportunity to increase training data by roughly 9.5$\times$. Second, the model was applied directly to the test set without any retraining or adaptation, meaning that any distributional shift between development and test data was unmitigated.


\subsection{Ablation Study}

To assess the contribution of individual components, we conducted an ablation study on the development set. Table~\ref{tab:ablation} reports accuracy when removing or replacing components from the full system. We note that all ablation results are on the development set and should be interpreted with caution given the observed overfitting.

\begin{table}[ht]
  \centering
  \begin{tabular}{lc}
    \hline
    \textbf{Configuration} & \textbf{Dev Acc} \\
    \hline
    Full system (Track A ensemble)        & $\sim$0.95  \\
    \quad $-$ KL distillation ($\lambda=0$) & $\sim$0.94  \\
    \quad $-$ Ensemble (single fold)       & $\sim$0.95  \\
    Track B only (cosine sim)              & $\sim$0.94  \\
    \quad $-$ Projection head (1024-dim)   & $\sim$0.93  \\
    \quad $-$ Composite loss (MNR only)    & $\sim$0.92  \\
    BGE-Large backbone (prev.\ system)     & $\sim$0.93  \\
    \hline
  \end{tabular}
  \caption{Ablation study on the development set. These results did not generalize to the test set.}
  \label{tab:ablation}
\end{table}

On the development set, the backbone upgrade from BGE-Large to Qwen3-Embedding-0.6B provided the single largest improvement, and the composite loss contributed meaningfully over using MNR alone. However, the test results suggest that these component-level improvements were largely artifacts of overfitting: the gains reflected memorization of development-set patterns rather than genuine improvements in narrative understanding.


\subsection{Error Analysis}

The near-chance test accuracy on Track~A (53.50\%) indicates that the system's learned representations do not generalize beyond the development distribution. We identify the following failure modes and contributing factors:

\paragraph{Overfitting to a small training set.} The system was trained on 200 labeled triples---still too few for a 0.6B-parameter encoder to learn robust narrative similarity patterns. The high development accuracy (Track~A CV mean 0.945; Track~B best 0.935) was misleading, as the model learned patterns that did not transfer to the hidden test distribution.

\paragraph{Unused synthetic training data.} The organizers provided 1,900 synthetic triples generated using large language models, which we did not incorporate. This data would have increased the training set by approximately 9.5$\times$, potentially providing the diversity and volume necessary to learn generalizable patterns. We consider this the single most impactful change we could have made.

\paragraph{No test-time adaptation.} The models trained on development data were applied directly to the test set without retraining or fine-tuning. Any distributional shift between the development and test narratives---in terms of genre, style, length, or thematic content---was therefore unmitigated.

\paragraph{Track B vs.\ Track A generalization.} Track~B (59.25\%) generalized somewhat better than Track~A (53.50\%), suggesting that the bi-encoder embeddings captured some useful narrative signal despite overfitting, while the MLP head introduced an additional layer of overfitting on top of already overfit embeddings.


\section{Conclusion}

We presented a two-stage system for SemEval-2026 Task 4: Narrative Similarity, built on Qwen3-Embedding-0.6B. The first stage fine-tunes the encoder as a bi-encoder with a 512-dimensional projection head using a composite loss (Track~B). The second stage adds a lightweight MLP head trained on frozen embeddings with pairwise interaction features and $k$-fold ensemble inference (Track~A). The system prioritizes parameter efficiency: the Track~A component contains fewer than 2 million parameters, and the Track~B projection head adds only $\sim$0.5 million parameters.

While the architecture achieved strong development accuracy ($\sim$0.95 on Track~A, $\sim$0.94 on Track~B), it generalized poorly to the official test set (53.50\% on Track~A, 59.25\% on Track~B). The primary cause was training on an insufficient amount of data: the system used only the small development set ($\sim$100 triples) and did not incorporate the 1,900 synthetic training triples provided by the organizers. This experience yields a clear lesson: for narrative similarity on limited data, training set size is more important than architectural sophistication.

Future work will focus on three directions: first, incorporating the synthetic training data and evaluating its impact on generalization; second, applying regularization techniques specifically designed for low-resource fine-tuning, such as R-Drop or mixup at the embedding level; and third, exploring whether larger Qwen3-Embedding variants (4B/8B) can better resist overfitting through stronger pretrained representations.


\section{Ethical Considerations}

Narrative similarity systems could potentially be misused for automated plagiarism of creative works, enabling large-scale identification of stylistically similar stories for appropriation rather than legitimate retrieval. Additionally, embedding-based similarity models may encode biases present in the training data: if certain narrative traditions, cultural contexts, or writing styles are underrepresented, the model may systematically underestimate the similarity between stories from those traditions.

Our system operates on short English-language narratives and has not been evaluated on other languages or longer literary works. The pretrained Qwen3-Embedding backbone was trained on a large multilingual corpus, and any biases present in that pretraining data may propagate to our fine-tuned model. Furthermore, the use of LLM-generated synthetic training data (which we did not use but was provided by the organizers) raises questions about the extent to which models trained on synthetic narratives may inherit the stylistic and cultural biases of the generating LLM. We encourage careful evaluation before deploying narrative similarity systems in contexts where biased similarity judgments could cause harm, such as automated content moderation or educational assessment.


\section{Acknowledgments}

% TODO: Add acknowledgments
The authors would like to acknowledge the support provided by the master's degree scholarship program in engineering at the Universidad Tecnologica de Bolivar (UTB) in Cartagena, Colombia.

% Custom bibliography entries only
\bibliography{references}

\appendix

%\section{Example Appendix}
%\label{sec:appendix}

%This is an appendix.

\end{document}
